\section{A joint orbit--index product law}

Fix real $\sigma>\sigma_0$ and write
\[
 S_{\Hspec}(\sigma)
 =\sum_\gamma e^{-\sigma\widehat\ell_\gamma}\Hspec_\gamma.
\]
This number is finite and positive.  Define
\begin{equation}
 \pi_\sigma(\gamma)
 =\frac{e^{-\sigma\widehat\ell_\gamma}\Hspec_\gamma}
 {S_{\Hspec}(\sigma)}.
 \label{eq:orbit-law}
\end{equation}
It is a canonical probability on the countable primitive orbit set.
Each primitive orbit appears once; marked points and repetitions are not
silently inserted into this marginal.

Equip the countable primitive-orbit set with the discrete topology.  Define
a probability on its product with $[0,\infty)$ by
\begin{equation}
 \nu_{\sigma,\tau}
 =\frac1{Z(\sigma,\tau)}
 \sum_\gamma\sum_{n\ge3}
 e^{-\sigma\widehat\ell_\gamma}b_{\gamma,n}e^{-\tau n}
 \delta_{(\gamma,\tau n)}.
 \label{eq:joint-measure}
\end{equation}

\begin{theorem}[Joint boundary profile]
As $\tau\downarrow0$,
\begin{equation}
 \nu_{\sigma,\tau}
 \Longrightarrow\pi_\sigma\otimes\Gamma(2,1),
 \label{eq:joint-limit}
\end{equation}
where $\Gamma(2,1)$ has density
$xe^{-x}\mathbf1_{x\ge0}\,\dd x$.
\end{theorem}

\begin{proof}
Let $F$ be bounded on the primitive orbit set and $r\ge0$.  Applying the
all-orbit theorem at $(1+r)\tau$ gives
\begin{align*}
 &\lim_{\tau\downarrow0}
 \int F(\gamma)e^{-rx}\,\dd\nu_{\sigma,\tau}(\gamma,x)\\
 &\qquad=\frac1{(1+r)^2}
 \sum_\gamma\pi_\sigma(\gamma)F(\gamma).
\end{align*}
The factor $(1+r)^{-2}$ is the Laplace transform of $\Gamma(2,1)$.

For completeness, tightness in the orbit coordinate follows from the
summable period majorant \eqref{eq:majorant}.  The crude packet estimate and
\[
 \tau^2\sum_{\tau n\ge R}ne^{-\tau n}\ll(R+1)e^{-R}
\]
give tightness in the scaled-index coordinate.  The displayed mixed
transforms therefore identify every subsequential limit, proving
\eqref{eq:joint-limit}.
\end{proof}

Two consequences are worth separating.  First, the scaled index has the same
$\Gamma(2,1)$ profile as the period-four packet, so the shape comes from the
totient main term rather than a special multiplier.  Second, the orbit
marginal is not the unweighted pressure distribution: it is tilted by the
complete Mahler spectral height.

The law \eqref{eq:orbit-law} remains trace-field arithmetic.  It does not
identify an orbit with a rational prime and does not preserve the individual
prime-ideal factors inside $D_{\gamma,n}$.
